% ==============================================================
% RELATÓRIO: 13 MODELOS DE REDES NEURAIS CONVOLUCIONAIS
% Tarefa de Classificação de Imagens de Moda (Bolsas x Calças)
% ==============================================================
\documentclass[12pt,a4paper]{article}

\usepackage[utf8]{inputenc}
\usepackage[T1]{fontenc}
\usepackage[brazil]{babel}
\usepackage{geometry}
\usepackage{graphicx}
\usepackage{booktabs}
\usepackage{longtable}
\usepackage{adjustbox}
\usepackage{caption}
\usepackage{subcaption}
\usepackage{hyperref}
\usepackage{amsmath}
\usepackage{amssymb}
\usepackage{float}
\usepackage{xcolor}
\usepackage{array}
\usepackage{multirow}
\usepackage{tabularx}
\usepackage{setspace}
\usepackage{fancyhdr}
\usepackage{titlesec}
\usepackage{parskip}
\usepackage{underscore}  % permite _ em modo texto sem escapar

\geometry{
  a4paper,
  left=2.5cm, right=2.5cm,
  top=3cm,    bottom=3cm
}

\graphicspath{%
  {../results/plots/}%
  {model_parameters/}%
}

\hypersetup{
  colorlinks=true,
  linkcolor=blue!70!black,
  urlcolor=blue!70!black,
  citecolor=blue!70!black,
  pdftitle={Relatório: 13 Modelos de Redes Neurais Convolucionais},
  pdfauthor={Anna Victoria}
}

\pagestyle{fancy}
\fancyhf{}
\rhead{\small Redes Neurais Convolucionais na Moda}
\lhead{\small Anna Victoria --- TCC}
\cfoot{\thepage}
\renewcommand{\headrulewidth}{0.4pt}

\titleformat{\section}{\large\bfseries\color{blue!70!black}}{}{0em}{\thesection\quad}
\titleformat{\subsection}{\normalsize\bfseries}{}{0em}{\thesubsection\quad}

% Comando auxiliar para inserir tabela de rodadas com ajuste de largura
\newcommand{\tabelarodadas}[2]{%
  \begin{table}[H]
    \centering
    \caption{#1}
    \begin{adjustbox}{max width=\textwidth, center}
      \input{#2}
    \end{adjustbox}
  \end{table}%
}

% Comando auxiliar para inserir tabela de resumo
\newcommand{\tabelaresumo}[2]{%
  \begin{table}[H]
    \centering
    \caption{#1}
    \begin{adjustbox}{max width=\textwidth, center}
      \input{#2}
    \end{adjustbox}
  \end{table}%
}

% Ambiente para bloco de parâmetros
\newenvironment{parametros}{%
  \begin{table}[H]
  \centering
  \small
  \begin{tabular}{ll}
  \toprule
  \textbf{Parâmetro} & \textbf{Valor} \\
  \midrule
}{%
  \bottomrule
  \end{tabular}
  \end{table}
}

\setlength{\parskip}{6pt}
\setlength{\parindent}{0pt}

% ==============================================================
\begin{document}
% ==============================================================

% ------ CAPA --------------------------------------------------
\begin{titlepage}
  \centering
  \vspace*{2cm}
  {\Huge\bfseries\color{blue!70!black}
    Redes Neurais Convolucionais\\[0.4em]
    na Classificação de Moda\\[0.4em]
    \large\normalfont\color{black}
    Relatório Técnico: Avaliação dos 13~Modelos\par}
  \vspace{1.5cm}
  {\large\bfseries Anna Victoria\par}
  \vspace{0.3cm}
  {\large Orientador: Fabrízzio Condé de Oliveira\par}
  \vspace{0.5cm}
  {\large Trabalho de Conclusão de Curso\par}
  \vspace{2cm}
  \begin{center}
    \begin{tabular}{lp{9cm}}
      \toprule
      \textbf{Tarefa} & Classificação binária: Bolsas \emph{vs.} Calças \\
      \textbf{Protocolo} & Repeated Holdout (10~sementes) \\
      \textbf{Modelos} & 13~(2~CNN do zero + 11~Transfer Learning) \\
      \textbf{Framework} & TensorFlow / Keras \\
      \textbf{Plataforma} & Google Colab (GPU) \\
      \bottomrule
    \end{tabular}
  \end{center}
  \vfill
  {\large\today\par}
\end{titlepage}

\tableofcontents
\clearpage

% ==============================================================
\section{Introdução}
% ==============================================================

Este relatório documenta a avaliação experimental de \textbf{13 modelos de Redes Neurais
Convolucionais (CNNs)} aplicados à classificação binária de imagens de moda, distinguindo
\textbf{bolsas} de \textbf{calças}. Os modelos compreendem duas categorias:

\begin{itemize}
  \item \textbf{CNNs treinadas do zero} (\emph{from scratch}): redes convolucionais projetadas
        e treinadas exclusivamente com os dados da tarefa, sem aproveitamento de pesos pré-treinados.
        Dois valores de taxa de aprendizado foram avaliados: $5\times10^{-4}$ e $1\times10^{-4}$.
  \item \textbf{Transfer Learning (TL)}: 11~modelos com \emph{backbones} pré-treinados na
        base ImageNet, cujas camadas convolucionais são congeladas e apenas a cabeça de
        classificação é treinada nos dados da tarefa.
\end{itemize}

O objetivo é identificar qual arquitetura oferece o melhor compromisso entre desempenho
preditivo (acurácia, precisão, recall e F1-score) e custo computacional (tempo de treinamento
e de inferência).

% ==============================================================
\section{Fundamentos Matemáticos}
\label{sec:fundamentos}
% ==============================================================

Esta seção apresenta as principais funções e expressões matemáticas que fundamentam
o funcionamento das CNNs avaliadas neste trabalho, desde a operação primitiva de
convolução até o otimizador e as funções estatísticas de replicabilidade.

% ------ Convolução 2D ------------------------------------------
\subsection{Operação de Convolução 2D}

A operação central de uma CNN é a \textbf{convolução discreta 2D}, que desliza um filtro
(kernel) $K$ de dimensão $k_h \times k_w$ sobre uma imagem de entrada $I$ e produz um
mapa de características:

\begin{equation}
  (I * K)[i,j] = \sum_{m=0}^{k_h-1}\sum_{n=0}^{k_w-1} I[i+m,\;j+n]\cdot K[m,n]
  \label{eq:conv2d}
\end{equation}

onde $i$ e $j$ indexam a posição espacial no mapa de saída. Em redes com múltiplos
canais de entrada $C$ e $N$ filtros paralelos, a equação se generaliza para:

\begin{equation}
  Z^{(n)}[i,j]
  = \sum_{c=1}^{C}\sum_{m=0}^{k_h-1}\sum_{n'=0}^{k_w-1}
    X^{(c)}[i+m,\;j+n']\cdot K^{(n,c)}[m,n'] + b^{(n)}
  \label{eq:conv2d_multi}
\end{equation}

onde $b^{(n)}$ é o viés do $n$-ésimo filtro. Neste projeto, todos os filtros
convolucionais utilizam núcleos $3\times3$.

% ------ Funções de Ativação ------------------------------------
\subsection{Funções de Ativação}

\subsubsection{ReLU --- \textit{Rectified Linear Unit}}

Após cada camada convolucional, aplica-se a função de ativação \textbf{ReLU},
que introduz não-linearidade ao modelo zerando todos os valores negativos:

\begin{equation}
  \text{ReLU}(x) = \max(0,\,x)
  \label{eq:relu}
\end{equation}

A ReLU é computacionalmente eficiente e mitiga o problema do gradiente desvanecente
em redes profundas, sendo a ativação padrão em todas as camadas ocultas dos 13 modelos.

\subsubsection{Logit --- Saída Linear Bruta da Rede}
\label{sec:logit}

Antes de aplicar qualquer função de ativação, a última camada densa da rede
(\texttt{Dense(2)}) computa uma combinação linear dos valores recebidos:

\begin{equation}
  z_c = \mathbf{w}_c^\top \mathbf{h} + b_c, \qquad c = 1, \ldots, C
  \label{eq:logit}
\end{equation}

onde $\mathbf{h}$ é o vetor de ativações da camada anterior (saída do
\texttt{Dense(64)}), $\mathbf{w}_c$ é o vetor de pesos associado à classe $c$
e $b_c$ é o viés correspondente.

O valor $z_c$ é chamado de \textbf{logit} da classe $c$. O termo vem da
função \emph{logit} da estatística ($\log\frac{p}{1-p}$), mas no contexto
de redes neurais designa genericamente qualquer \textbf{pontuação linear
irrestrita} produzida antes de uma função de ativação --- podendo assumir
qualquer valor real ($z_c \in \mathbb{R}$), positivo ou negativo, sem
interpretação direta de probabilidade.

O vetor de logits $\mathbf{z} = (z_1, z_2)$ para a tarefa binária
deste trabalho representa, portanto, o \emph{grau de afinidade} bruto
que a rede atribui a cada classe. Um logit maior indica maior tendência
da rede em classificar a entrada naquela classe, mas apenas após a
Softmax esses valores se tornam probabilidades comparáveis e somam 1.

\subsubsection{Softmax --- Conversão para Probabilidade}

A camada de saída de \textbf{todos os 13 modelos} utiliza a função \textbf{Softmax},
que converte o vetor de logits $\mathbf{z} = (z_1, z_2, \ldots, z_C)$ em uma
distribuição de probabilidade sobre as $C$ classes:

\begin{equation}
  \hat{p}_c = \operatorname{softmax}(z_c)
            = \frac{e^{z_c}}{\displaystyle\sum_{j=1}^{C} e^{z_j}},
  \qquad c = 1, \ldots, C
  \label{eq:softmax}
\end{equation}

A Softmax garante que $\hat{p}_c \in (0,1)$ para todo $c$ e que
$\sum_{c=1}^{C}\hat{p}_c = 1$. Para a tarefa binária deste trabalho
($C = 2$: bolsas e calças), as duas probabilidades complementares são:

\begin{equation}
  \hat{p}_{\text{bolsa}} = \frac{e^{z_1}}{e^{z_1}+e^{z_2}},
  \qquad
  \hat{p}_{\text{calça}} = \frac{e^{z_2}}{e^{z_1}+e^{z_2}}
  \label{eq:softmax_bin}
\end{equation}

A classe predita é determinada por $\hat{y} = \arg\max_{c}\,\hat{p}_c$.

% ------ Pooling ------------------------------------------------
\subsection{Operações de Pooling}

\subsubsection{Max Pooling}

O \textbf{Max Pooling} reduz a dimensão espacial do mapa de características
retendo o valor máximo em cada janela $p\times p$ (neste projeto $p=2$) com
passo (\emph{stride}) igual a $p$:

\begin{equation}
  y[i,j] = \max_{\substack{0\le m<p\\0\le n<p}}\; x[p\,i+m,\;p\,j+n]
  \label{eq:maxpool}
\end{equation}

Cada bloco Conv+MaxPool na CNN do zero reduz pela metade a resolução espacial
($H\times W\to\tfrac{H}{2}\times\tfrac{W}{2}$), expandindo progressivamente
o campo receptivo da rede.

\subsubsection{Global Average Pooling (GAP)}

Nos modelos de Transfer Learning, o backbone produz um tensor
$F\in\mathbb{R}^{H\times W\times C}$. A camada \texttt{GlobalAveragePooling2D}
colapsa as dimensões espaciais calculando a média por canal, gerando um
vetor $\mathbf{g}\in\mathbb{R}^{C}$:

\begin{equation}
  g_c = \frac{1}{H\cdot W}\sum_{i=1}^{H}\sum_{j=1}^{W} F[i,j,c],
  \qquad c = 1,\ldots,C
  \label{eq:gap}
\end{equation}

O GAP elimina a necessidade de camadas densas volumosas sobre o mapa de
características e atua como regularizador natural, reduzindo o risco de
\emph{overfitting}.

% ------ Função de Perda ----------------------------------------
\subsection{Função de Perda: Entropia Cruzada Categórica}

Todos os modelos são treinados minimizando a \textbf{Entropia Cruzada
Categórica} (\emph{Categorical Cross-Entropy}), que mede a divergência
entre a distribuição real $\mathbf{y}$ (codificada em \emph{one-hot})
e a distribuição prevista $\hat{\mathbf{p}}$ (saída da Softmax):

\begin{equation}
  \mathcal{L} = -\frac{1}{N}\sum_{i=1}^{N}\sum_{c=1}^{C}
                y_{i,c}\,\log\!\left(\hat{p}_{i,c}\right)
  \label{eq:crossentropy}
\end{equation}

onde $N$ é o número de amostras no mini-batch, $y_{i,c}\in\{0,1\}$ indica se
a amostra $i$ pertence à classe $c$, e $\hat{p}_{i,c}$ é a probabilidade
prevista pela Softmax. Como $\mathbf{y}$ é \emph{one-hot}, apenas o termo
da classe verdadeira $c_i^*$ contribui para a soma:

\begin{equation}
  \mathcal{L} = -\frac{1}{N}\sum_{i=1}^{N}\log\!\left(\hat{p}_{i,\,c_i^*}\right)
  \label{eq:crossentropy_simp}
\end{equation}

A perda é minimizada quando $\hat{p}_{i,c_i^*}\to 1$, isto é, quando o modelo
atribui probabilidade máxima à classe correta. O Early Stopping monitora
$\mathcal{L}$ na faixa de treino com \texttt{patience}~$= 5$.

% ------ Exemplo Numérico ---------------------------------------
\subsection{Exemplo Numérico: Softmax, Probabilidade e Função de Perda}
\label{sec:exemplo_numerico}

Para concretizar a relação entre logits, Softmax e Entropia Cruzada, considere uma
imagem de \textbf{bolsa} (classe verdadeira $c^* = \text{bolsa}$, codificação
\emph{one-hot} $\mathbf{y} = [1,\;0]$) passando pela camada de saída
\texttt{Dense(2, activation='softmax')}, que produz dois logits
$\mathbf{z} = (z_{\text{bolsa}},\; z_{\text{calça}})$.
Analisamos três cenários que ilustram como a Softmax converte logits em
probabilidades e como a Entropia Cruzada penaliza cada situação.

\vspace{0.8em}
\noindent\textbf{Cenário~1 --- Predição correta (modelo confiante):}
\[
  \mathbf{z} = (2{,}10;\; 0{,}30)
\]
\begin{align*}
  \hat{p}_{\text{bolsa}}
    &= \frac{e^{2{,}10}}{e^{2{,}10}+e^{0{,}30}}
     = \frac{8{,}166}{8{,}166 + 1{,}350}
     = \frac{8{,}166}{9{,}516}
     \approx 0{,}858 \\[6pt]
  \hat{p}_{\text{calça}}
    &= \frac{e^{0{,}30}}{9{,}516}
     = \frac{1{,}350}{9{,}516}
     \approx 0{,}142
\end{align*}
\[
  \mathcal{L} = -\log(\hat{p}_{c^*}) = -\log(0{,}858) \approx \mathbf{0{,}153}
\]
A perda é baixa: o modelo atribui 85,8\,\% de probabilidade à classe
correta. O gradiente resultante é pequeno e o Adam realiza uma atualização
sutil dos pesos.

\vspace{0.8em}
\noindent\textbf{Cenário~2 --- Predição incerta (modelo sem confiança):}
\[
  \mathbf{z} = (0{,}50;\; 0{,}50)
\]
\begin{align*}
  \hat{p}_{\text{bolsa}}
    &= \frac{e^{0{,}50}}{e^{0{,}50}+e^{0{,}50}}
     = \frac{1{,}649}{3{,}297}
     = 0{,}500 \\[6pt]
  \hat{p}_{\text{calça}}
    &= \frac{1{,}649}{3{,}297}
     = 0{,}500
\end{align*}
\[
  \mathcal{L} = -\log(0{,}500) \approx \mathbf{0{,}693}
\]
Logits iguais produzem distribuição uniforme --- máxima incerteza.
A perda moderada indica que o modelo ainda não aprendeu a distinguir
as classes, situação típica das primeiras épocas de treinamento.

\vspace{0.8em}
\noindent\textbf{Cenário~3 --- Predição errada (modelo confiante na classe errada):}
\[
  \mathbf{z} = (-1{,}20;\; 2{,}50)
\]
\begin{align*}
  \hat{p}_{\text{bolsa}}
    &= \frac{e^{-1{,}20}}{e^{-1{,}20}+e^{2{,}50}}
     = \frac{0{,}301}{0{,}301 + 12{,}182}
     = \frac{0{,}301}{12{,}483}
     \approx 0{,}024 \\[6pt]
  \hat{p}_{\text{calça}}
    &= \frac{12{,}182}{12{,}483}
     \approx 0{,}976
\end{align*}
\[
  \mathcal{L} = -\log(0{,}024) \approx \mathbf{3{,}730}
\]
A perda é alta: apenas 2,4\,\% de probabilidade foi atribuída à classe
correta. O gradiente é grande, forçando o Adam a realizar uma correção
expressiva nos pesos $\theta$ da rede.

\vspace{1em}
A Tabela~\ref{tab:exemplo_loss} resume os três cenários e evidencia a
relação monotônica entre \emph{confiança na classe errada} e \emph{magnitude da perda}:

\begin{table}[H]
\centering
\caption{Resumo do exemplo numérico --- logits, probabilidades Softmax e Entropia
         Cruzada para uma imagem de bolsa ($\mathbf{y}=[1,0]$)}
\label{tab:exemplo_loss}
\begin{tabular}{lcccccc}
\toprule
\textbf{Cenário}
  & $z_{\text{bolsa}}$ & $z_{\text{calça}}$
  & $\hat{p}_{\text{bolsa}}$ & $\hat{p}_{\text{calça}}$
  & $\mathcal{L}$ & \textbf{Situação} \\
\midrule
Predição correta  & $2{,}10$  & $0{,}30$ & $0{,}858$ & $0{,}142$ & $0{,}153$ & \checkmark \\
Predição incerta  & $0{,}50$  & $0{,}50$ & $0{,}500$ & $0{,}500$ & $0{,}693$ & $\sim$     \\
Predição errada   & $-1{,}20$ & $2{,}50$ & $0{,}024$ & $0{,}976$ & $3{,}730$ & $\times$   \\
\bottomrule
\end{tabular}
\end{table}

\noindent\textbf{Interpretação geral:} à medida que o treinamento progride,
o Adam minimiza $\mathcal{L}$ ajustando os pesos para que o logit da classe
correta cresça e o da classe errada diminua, empurrando
$\hat{p}_{c^*}\to 1$ e $\mathcal{L}\to 0$.
Para o modelo \emph{perfeito} ($\hat{p}_{c^*}=1$),
a perda seria $-\log(1) = 0$; para o modelo totalmente errado
($\hat{p}_{c^*}\to 0$), a perda tenderia a $+\infty$.

% ------ Adam ---------------------------------------------------
\subsection{Otimizador Adam}

O treinamento utiliza o otimizador \textbf{Adam}
(\emph{Adaptive Moment Estimation}) \cite{kingma2014adam}, que adapta
individualmente a taxa de aprendizado de cada parâmetro com base nos
primeiros e segundos momentos dos gradientes.

Seja $g_t = \nabla_\theta\mathcal{L}_t$ o gradiente no passo $t$.
As equações de atualização são:

\begin{align}
  m_t       &= \beta_1\,m_{t-1} + (1-\beta_1)\,g_t
                                                           \label{eq:adam_m}\\[3pt]
  v_t       &= \beta_2\,v_{t-1} + (1-\beta_2)\,g_t^2
                                                           \label{eq:adam_v}\\[3pt]
  \hat{m}_t &= \dfrac{m_t}{1-\beta_1^t}
                                                           \label{eq:adam_mhat}\\[3pt]
  \hat{v}_t &= \dfrac{v_t}{1-\beta_2^t}
                                                           \label{eq:adam_vhat}\\[6pt]
  \theta_t  &= \theta_{t-1}
               - \alpha\cdot\dfrac{\hat{m}_t}{\sqrt{\hat{v}_t}+\varepsilon}
                                                           \label{eq:adam_update}
\end{align}

onde $m_t$ é o \textbf{primeiro momento} (média exponencial dos gradientes),
$v_t$ é o \textbf{segundo momento} (variância exponencial), e as correções
$\hat{m}_t$, $\hat{v}_t$ compensam o viés de inicialização em zero.
Os hiperparâmetros padrão do Keras são
$\beta_1 = 0{,}9$, $\beta_2 = 0{,}999$, $\varepsilon = 10^{-7}$, e $\alpha$
é a taxa de aprendizado ($10^{-4}$ para os modelos de TL;
$5\times10^{-4}$ ou $10^{-4}$ para as CNNs do zero).

% ------ Convolução Separável -----------------------------------
\subsection{Convolução Separável em Profundidade}

Os modelos da família MobileNet (V2, V3Small e V3Large) e o Xception
substituem as convoluções padrão por \textbf{convoluções separáveis em
profundidade} (\emph{Depthwise Separable Convolutions})
\cite{chollet2017xception}, compostas de dois passos sequenciais:

\textbf{(1) Depthwise Convolution} --- aplica um filtro independente por canal de entrada:
\begin{equation}
  \hat{Z}^{(c)}[i,j]
  = \sum_{m}\sum_{n} X^{(c)}[i+m,\,j+n]\cdot K_{\mathrm{dw}}^{(c)}[m,n]
  \label{eq:depthwise}
\end{equation}

\textbf{(2) Pointwise Convolution} --- combina os canais com uma convolução $1\times1$:
\begin{equation}
  Z^{(n)}[i,j] = \sum_{c=1}^{C}\hat{Z}^{(c)}[i,j]\cdot K_{\mathrm{pw}}^{(n,c)}
  \label{eq:pointwise}
\end{equation}

Sendo $D_K$ o tamanho do kernel, $C$ os canais de entrada, $N$ o número de
filtros e $D_F\times D_F$ a resolução do mapa, a razão de custo
computacional entre as duas abordagens é:

\begin{equation}
  \frac{\mathcal{C}_{\mathrm{sep}}}{\mathcal{C}_{\mathrm{padrão}}}
  = \frac{D_K^2 C D_F^2 + C N D_F^2}{D_K^2 C N D_F^2}
  = \frac{1}{N} + \frac{1}{D_K^2}
  \label{eq:dw_ratio}
\end{equation}

Para $D_K = 3$ e $N \gg 1$, a redução é de aproximadamente $\mathbf{8\times}$,
viabilizando o uso em dispositivos com restrições de memória e processamento.

% ------ Batch Normalization ------------------------------------
\subsection{Normalização em Lote (\textit{Batch Normalization})}

Vários backbones (ResNet50, EfficientNet, DenseNet121 e outros) aplicam
\textbf{Batch Normalization} após as convoluções, normalizando as ativações
dentro do mini-batch para reduzir o \emph{internal covariate shift}:

\begin{equation}
  \hat{x}_i = \frac{x_i - \mu_B}{\sqrt{\sigma_B^2 + \varepsilon}}
  \label{eq:bn_norm}
\end{equation}

onde $\mu_B$ e $\sigma_B^2$ são a média e a variância das ativações no
batch atual. Em seguida, uma transformação afim com parâmetros aprendíveis
$\gamma$ (escala) e $\beta$ (deslocamento) restaura a capacidade expressiva:

\begin{equation}
  y_i = \gamma\,\hat{x}_i + \beta
  \label{eq:bn_affine}
\end{equation}

% ------ Estatísticas de Replicabilidade ------------------------
\subsection{Estatísticas de Replicabilidade}

O protocolo Repeated Holdout com $n = 10$ rodadas exige a quantificação da
variabilidade dos resultados. As seguintes medidas são reportadas para cada métrica:

\textbf{Desvio Padrão Amostral:}
\begin{equation}
  s = \sqrt{\frac{1}{n-1}\sum_{i=1}^{n}(x_i - \bar{x})^2}
  \label{eq:std}
\end{equation}

\textbf{Coeficiente de Variação (CV):}
\begin{equation}
  \mathrm{CV} = \frac{s}{\bar{x}}\times 100\%
  \label{eq:cv}
\end{equation}

O CV expressa a dispersão relativa à média, permitindo comparar a estabilidade
entre modelos com desempenhos médios distintos. Valores baixos de CV indicam
maior estabilidade frente à variação das sementes.

\textbf{Intervalo de Confiança de 95\% (t-Student):}
\begin{equation}
  \mathrm{IC}_{95\%} = \bar{x} \pm t_{0{,}025,\;n-1}\cdot\frac{s}{\sqrt{n}}
  \label{eq:ic95_fund}
\end{equation}

Para $n = 10$ rodadas utiliza-se $t_{0{,}025,\;9} = 2{,}262$.

% ==============================================================
\section{Metodologia Experimental}
% ==============================================================

\subsection{Protocolo de Avaliação: Repeated Holdout}

O protocolo \textbf{Repeated Holdout} com 10~sementes aleatórias fixas
($\text{SEED} \in \{0, 1, \ldots, 9\}$) foi adotado para garantir a
\textbf{replicabilidade} e estimar a variabilidade dos resultados.
Em cada rodada, uma semente diferente é utilizada para inicializar
os geradores aleatórios do Python, NumPy e TensorFlow, assegurando
que os resultados sejam reprodutíveis.

\subsection{Conjunto de Dados e Divisão}

\begin{itemize}
  \item \textbf{Treinamento}: imagens nas pastas \texttt{BOLSAS TREINAMENTO} e
        \texttt{CALÇAS TREINAMENTO}, localizadas em \texttt{TREINAMENTO/}.
  \item \textbf{Teste}: imagens nas pastas \texttt{BOLSAS TESTE} e
        \texttt{CALÇAS TESTE}, localizadas em \texttt{TESTE/}.
  \item A divisão treino/teste é \emph{fixa}; apenas a semente de embaralhamento
        do gerador de treino varia entre rodadas.
\end{itemize}

\subsection{Aumento de Dados (\textit{Data Augmentation})}

Dois conjuntos de parâmetros de aumento foram utilizados:

\begin{table}[H]
\centering
\caption{Parâmetros de aumento de dados por tipo de modelo}
\begin{tabular}{lcc}
\toprule
\textbf{Parâmetro} & \textbf{CNN do Zero} & \textbf{Transfer Learning} \\
\midrule
\texttt{rotation\_range}     & 15 & 15 \\
\texttt{zoom\_range}         & 0.25 & 0.20 \\
\texttt{width\_shift\_range} & 0.10 & 0.10 \\
\texttt{height\_shift\_range}& 0.10 & 0.10 \\
\texttt{shear\_range}        & 0.20 & --- \\
\texttt{horizontal\_flip}    & \checkmark & \checkmark \\
\texttt{brightness\_range}   & [0.8, 1.2] & --- \\
\bottomrule
\end{tabular}
\end{table}

\subsection{Pré-processamento de Imagens}

Cada família de backbone exige uma função de pré-processamento específica,
que normaliza os pixels para a escala esperada durante o pré-treino no ImageNet.
As CNNs do zero utilizam rescaling simples ($\div 255$).

\begin{table}[H]
\centering
\caption{Funções de pré-processamento por backbone}
\begin{tabular}{ll}
\toprule
\textbf{Backbone} & \textbf{Função de pré-processamento} \\
\midrule
CNN do zero            & \texttt{rescale = 1/255} \\
MobileNetV2            & \texttt{mobilenet\_v2.preprocess\_input} \\
EfficientNetB0/B1/B2   & \texttt{efficientnet.preprocess\_input} \\
ResNet50               & \texttt{resnet.preprocess\_input} \\
DenseNet121            & \texttt{densenet.preprocess\_input} \\
MobileNetV3Small/Large & \texttt{mobilenet\_v3.preprocess\_input} \\
NASNetMobile           & \texttt{nasnet.preprocess\_input} \\
InceptionV3            & \texttt{inception\_v3.preprocess\_input} \\
Xception               & \texttt{xception.preprocess\_input} \\
\bottomrule
\end{tabular}
\end{table}

\subsection{Callback: Early Stopping}

Todos os modelos utilizam \texttt{EarlyStopping} com os seguintes parâmetros:
\begin{itemize}
  \item \texttt{monitor = 'loss'} (monitora a perda de treino)
  \item \texttt{patience = 5} (aguarda 5 épocas sem melhoria antes de parar)
  \item \texttt{restore\_best\_weights = True}
\end{itemize}

\subsection{Métricas de Avaliação}

Quatro métricas de classificação são calculadas com média ponderada (\emph{weighted average}):
\begin{align}
  \text{Acurácia}  &= \frac{TP + TN}{TP + TN + FP + FN} \\[4pt]
  \text{Precisão}  &= \frac{TP}{TP + FP} \\[4pt]
  \text{Recall}    &= \frac{TP}{TP + FN} \\[4pt]
  \text{F1-score}  &= \frac{2 \cdot \text{Precisão} \cdot \text{Recall}}{\text{Precisão} + \text{Recall}}
\end{align}

Além disso, são registrados os tempos de treinamento, inferência e total (em segundos).

\subsection{Intervalo de Confiança}

O intervalo de confiança de 95\% é calculado usando a distribuição \emph{t}-Student
para amostras pequenas ($n = 10$), com $t_{0.025, 9} = 2{,}262$:
\[
  \text{IC}_{95\%} = \bar{x} \pm t_{0.025,\,n-1} \cdot \dfrac{s}{\sqrt{n}}
\]

\subsection{Configuração Geral dos 13 Modelos}

\begin{table}[H]
\centering
\caption{Resumo de configuração de todos os modelos avaliados}
\small
\begin{tabular}{clllccccl}
\toprule
\textbf{ID} & \textbf{Modelo} & \textbf{Tipo} & \textbf{Backbone} &
\textbf{Img} & \textbf{Batch} & \textbf{Ép.} & \textbf{LR} & \textbf{Preprocess} \\
\midrule
1  & cnn\_from\_scratch      & CNN & ConvNet        & 150 & 32 & 50 & $5\times10^{-4}$ & Rescale  \\
2  & cnn\_from\_scratch\_lr1e4 & CNN & ConvNet      & 150 & 32 & 50 & $1\times10^{-4}$ & Rescale  \\
3  & mobilenetv2             & TL  & MobileNetV2    & 150 & 32 & 30 & $1\times10^{-4}$ & MNv2     \\
4  & efficientnetb0          & TL  & EfficientNetB0 & 224 & 32 & 30 & $1\times10^{-4}$ & EffNet   \\
5  & efficientnetb1          & TL  & EfficientNetB1 & 240 & 32 & 30 & $1\times10^{-4}$ & EffNet   \\
6  & efficientnetb2          & TL  & EfficientNetB2 & 260 & 32 & 30 & $1\times10^{-4}$ & EffNet   \\
7  & resnet50                & TL  & ResNet50       & 224 & 32 & 30 & $1\times10^{-4}$ & ResNet   \\
8  & densenet121             & TL  & DenseNet121    & 224 & 32 & 30 & $1\times10^{-4}$ & DenseNet \\
9  & mobilenetv3small        & TL  & MobileNetV3S   & 224 & 32 & 30 & $1\times10^{-4}$ & MNv3     \\
10 & mobilenetv3large        & TL  & MobileNetV3L   & 224 & 32 & 30 & $1\times10^{-4}$ & MNv3     \\
11 & nasnetmobile            & TL  & NASNetMobile   & 224 & 32 & 30 & $1\times10^{-4}$ & NASNet   \\
12 & inceptionv3             & TL  & InceptionV3    & 299 & 16 & 30 & $1\times10^{-4}$ & Inception\\
13 & xception                & TL  & Xception       & 299 & 16 & 30 & $1\times10^{-4}$ & Xception \\
\bottomrule
\end{tabular}
\end{table}

\subsection{Arquitetura da Cabeça de Classificação (Transfer Learning)}

Para todos os modelos de Transfer Learning, a mesma cabeça de classificação
é adicionada sobre o backbone congelado:

\begin{enumerate}
  \item \textbf{Backbone pré-treinado} (pesos ImageNet, \texttt{trainable = False},
        \texttt{include\_top = False})
  \item \texttt{GlobalAveragePooling2D()} --- reduz o mapa de características a um vetor
  \item \texttt{Dense(128, activation='relu')}
  \item \texttt{Dense(64, activation='relu')}
  \item \texttt{Dense(2, activation='softmax')} --- saída binária
\end{enumerate}

Compilação: \texttt{Adam(lr=1e-4)}, \texttt{loss='categorical\_crossentropy'},
\texttt{metrics=['accuracy']}.

% ==============================================================
\section{Modelo 1 --- CNN Treinada do Zero (lr $= 5\times10^{-4}$)}
\label{sec:cnn_scratch_5e4}
% ==============================================================

\subsection{Descrição da Arquitetura}

A CNN do zero (\texttt{cnn\_from\_scratch}) é uma rede convolucional simples,
projetada sem aproveitar pesos pré-treinados. Sua arquitetura segue o padrão
clássico de empilhamento de blocos Conv+Pool, finalizando com camadas totalmente
conectadas:

\begin{enumerate}
  \item \texttt{Conv2D(32, (3,3), activation='relu')} + \texttt{MaxPool2D()}
  \item \texttt{Conv2D(64, (3,3), activation='relu')} + \texttt{MaxPool2D()}
  \item \texttt{Conv2D(128, (3,3), activation='relu')} + \texttt{MaxPool2D()}
  \item \texttt{Flatten()}
  \item \texttt{Dense(128, activation='relu')}
  \item \texttt{Dense(64, activation='relu')}
  \item \texttt{Dense(2, activation='softmax')}
\end{enumerate}

\subsection{Hiperparâmetros}

\begin{parametros}
Taxa de Aprendizado   & $5 \times 10^{-4}$ \\
Imagem de Entrada     & $150 \times 150 \times 3$ \\
Tamanho do Batch      & 32 \\
Épocas Máximas        & 50 \\
Otimizador            & Adam \\
Função de Perda       & Categorical Crossentropy \\
Pré-processamento     & Rescale $1/255$ \\
\end{parametros}

\subsection{Resultados por Rodada}

\tabelarodadas{Resultados por rodada --- CNN do zero ($\text{lr}=5\times10^{-4}$)}%
              {../results/latex/tabela_rodadas_cnn_from_scratch.tex}

\subsection{Estatísticas Resumo}

\tabelaresumo{Resumo estatístico --- CNN do zero ($\text{lr}=5\times10^{-4}$)}%
             {../results/latex/tabela_resumo_cnn_from_scratch.tex}

\subsection{Gráficos}

\begin{figure}[H]
  \centering
  \includegraphics[width=0.85\textwidth]{cnn_from_scratch_01_acuracia_por_rodada.pdf}
  \caption{Acurácia por rodada --- CNN do zero ($\text{lr}=5\times10^{-4}$)}
\end{figure}

\begin{figure}[H]
  \centering
  \includegraphics[width=0.85\textwidth]{cnn_from_scratch_02_boxplot_metricas.pdf}
  \caption{Boxplot das métricas --- CNN do zero ($\text{lr}=5\times10^{-4}$)}
\end{figure}

\begin{figure}[H]
  \centering
  \includegraphics[width=0.85\textwidth]{cnn_from_scratch_03_media_mais_desvio.pdf}
  \caption{Média $\pm$ desvio padrão --- CNN do zero ($\text{lr}=5\times10^{-4}$)}
\end{figure}

\subsection{Análise}

A CNN treinada do zero com $\text{lr}=5\times10^{-4}$ apresentou
\textbf{alta variabilidade} entre as rodadas, com acurácia média de $0{,}7400$
e desvio padrão de $0{,}0903$ (CV = 12,2\%). O intervalo de confiança de 95\%
para a acurácia é $[0{,}6754;\, 0{,}8046]$. Esse comportamento indica que o
modelo é sensível à inicialização dos pesos e à seleção dos dados de treinamento,
resultado esperado para redes sem pré-treinamento. A taxa de aprendizado mais
alta ($5\times10^{-4}$) pode causar oscilações na otimização.

% ==============================================================
\section{Modelo 2 --- CNN Treinada do Zero (lr $= 1\times10^{-4}$)}
\label{sec:cnn_scratch_1e4}
% ==============================================================

\subsection{Descrição da Arquitetura}

Arquitetura idêntica ao Modelo~1 (seção~\ref{sec:cnn_scratch_5e4}),
diferindo apenas na taxa de aprendizado, que é reduzida para $1\times10^{-4}$.
O objetivo é verificar se um gradiente mais conservador melhora a estabilidade
do treinamento.

\subsection{Hiperparâmetros}

\begin{parametros}
Taxa de Aprendizado   & $1 \times 10^{-4}$ \\
Imagem de Entrada     & $150 \times 150 \times 3$ \\
Tamanho do Batch      & 32 \\
Épocas Máximas        & 50 \\
Otimizador            & Adam \\
Função de Perda       & Categorical Crossentropy \\
Pré-processamento     & Rescale $1/255$ \\
\end{parametros}

\subsection{Resultados por Rodada}

\tabelarodadas{Resultados por rodada --- CNN do zero ($\text{lr}=1\times10^{-4}$)}%
              {../results/latex/tabela_rodadas_cnn_from_scratch_lr1e4.tex}

\subsection{Estatísticas Resumo}

\tabelaresumo{Resumo estatístico --- CNN do zero ($\text{lr}=1\times10^{-4}$)}%
             {../results/latex/tabela_resumo_cnn_from_scratch_lr1e4.tex}

\subsection{Gráficos}

\begin{figure}[H]
  \centering
  \includegraphics[width=0.85\textwidth]{cnn_from_scratch_lr1e4_01_acuracia_por_rodada.pdf}
  \caption{Acurácia por rodada --- CNN do zero ($\text{lr}=1\times10^{-4}$)}
\end{figure}

\begin{figure}[H]
  \centering
  \includegraphics[width=0.85\textwidth]{cnn_from_scratch_lr1e4_02_boxplot_metricas.pdf}
  \caption{Boxplot das métricas --- CNN do zero ($\text{lr}=1\times10^{-4}$)}
\end{figure}

\begin{figure}[H]
  \centering
  \includegraphics[width=0.85\textwidth]{cnn_from_scratch_lr1e4_03_media_mais_desvio.pdf}
  \caption{Média $\pm$ desvio padrão --- CNN do zero ($\text{lr}=1\times10^{-4}$)}
\end{figure}

\subsection{Análise}

Com $\text{lr}=1\times10^{-4}$, a CNN do zero obteve acurácia média de $0{,}6567$,
\textbf{inferior} ao resultado com lr$=5\times10^{-4}$ ($0{,}7400$). O coeficiente
de variação foi de 12,9\%, similar ao Modelo~1. Isso sugere que a taxa de aprendizado
mais baixa pode ter dificultado a convergência dentro das 50~épocas com early stopping,
ou levado o modelo a mínimos locais subótimos. Em ambos os casos de CNN do zero,
o desempenho fica muito abaixo dos modelos de Transfer Learning.

% ==============================================================
\section{Modelo 3 --- MobileNetV2}
\label{sec:mobilenetv2}
% ==============================================================

\subsection{Descrição do Backbone}

O \textbf{MobileNetV2} \cite{sandler2018mobilenetv2} é uma arquitetura eficiente para dispositivos
móveis, baseada em \emph{inverted residuals} com conexões residuais em blocos de expansão linear.
O operador principal é o \emph{depthwise separable convolution}, que reduz drasticamente o número
de parâmetros em relação a convoluções padrão. A versão utilizada é a original com entrada
de $150\times150$.

\subsection{Hiperparâmetros}

\begin{parametros}
Backbone              & MobileNetV2 (ImageNet, congelado) \\
Taxa de Aprendizado   & $1 \times 10^{-4}$ \\
Imagem de Entrada     & $150 \times 150 \times 3$ \\
Tamanho do Batch      & 32 \\
Épocas Máximas        & 30 \\
Otimizador            & Adam \\
Pré-processamento     & \texttt{mobilenet\_v2.preprocess\_input} \\
\end{parametros}

\subsection{Resultados por Rodada}

\tabelarodadas{Resultados por rodada --- MobileNetV2}%
              {../results/latex/tabela_rodadas_mobilenetv2.tex}

\subsection{Estatísticas Resumo}

\tabelaresumo{Resumo estatístico --- MobileNetV2}%
             {../results/latex/tabela_resumo_mobilenetv2.tex}

\subsection{Gráficos}

\begin{figure}[H]
  \centering
  \includegraphics[width=0.85\textwidth]{mobilenetv2_01_acuracia_por_rodada.pdf}
  \caption{Acurácia por rodada --- MobileNetV2}
\end{figure}

\begin{figure}[H]
  \centering
  \includegraphics[width=0.85\textwidth]{mobilenetv2_02_boxplot_metricas.pdf}
  \caption{Boxplot das métricas --- MobileNetV2}
\end{figure}

\begin{figure}[H]
  \centering
  \includegraphics[width=0.85\textwidth]{mobilenetv2_03_media_mais_desvio.pdf}
  \caption{Média $\pm$ desvio padrão --- MobileNetV2}
\end{figure}

\subsection{Análise}

O MobileNetV2 atingiu acurácia média de $0{,}9700$ com desvio padrão de $0{,}0205$
(CV = 2,1\%), demonstrando estabilidade significativamente superior às CNNs do zero.
O intervalo de confiança de 95\% é $[0{,}9553;\, 0{,}9847]$. O tempo médio de
treinamento foi de $40{,}3$~s, tornando-o um dos modelos mais rápidos entre os de TL.

% ==============================================================
\section{Modelo 4 --- EfficientNetB0}
\label{sec:efficientnetb0}
% ==============================================================

\subsection{Descrição do Backbone}

O \textbf{EfficientNet} \cite{tan2019efficientnet} introduz o conceito de
\emph{compound scaling}, escalando simultaneamente profundidade, largura e
resolução da rede de forma balanceada. O \textbf{EfficientNetB0} é a versão
base desta família, com resolução nativa de $224\times224$.

\subsection{Hiperparâmetros}

\begin{parametros}
Backbone              & EfficientNetB0 (ImageNet, congelado) \\
Taxa de Aprendizado   & $1 \times 10^{-4}$ \\
Imagem de Entrada     & $224 \times 224 \times 3$ \\
Tamanho do Batch      & 32 \\
Épocas Máximas        & 30 \\
Otimizador            & Adam \\
Pré-processamento     & \texttt{efficientnet.preprocess\_input} \\
\end{parametros}

\subsection{Resultados por Rodada}

\tabelarodadas{Resultados por rodada --- EfficientNetB0}%
              {../results/latex/tabela_rodadas_efficientnetb0.tex}

\subsection{Estatísticas Resumo}

\tabelaresumo{Resumo estatístico --- EfficientNetB0}%
             {../results/latex/tabela_resumo_efficientnetb0.tex}

\subsection{Gráficos}

\begin{figure}[H]
  \centering
  \includegraphics[width=0.85\textwidth]{efficientnetb0_01_acuracia_por_rodada.pdf}
  \caption{Acurácia por rodada --- EfficientNetB0}
\end{figure}

\begin{figure}[H]
  \centering
  \includegraphics[width=0.85\textwidth]{efficientnetb0_02_boxplot_metricas.pdf}
  \caption{Boxplot das métricas --- EfficientNetB0}
\end{figure}

\begin{figure}[H]
  \centering
  \includegraphics[width=0.85\textwidth]{efficientnetb0_03_media_mais_desvio.pdf}
  \caption{Média $\pm$ desvio padrão --- EfficientNetB0}
\end{figure}

\subsection{Análise}

O EfficientNetB0 obteve acurácia média de $0{,}9983$, com desvio padrão de apenas
$0{,}0053$ (CV = 0,5\%), mostrando desempenho quase perfeito e alta estabilidade.
O IC~95\% é $[0{,}9946;\, 1{,}0021]$. O tempo médio de treinamento foi de $47{,}1$~s,
levemente superior ao MobileNetV2.

% ==============================================================
\section{Modelo 5 --- EfficientNetB1}
\label{sec:efficientnetb1}
% ==============================================================

\subsection{Descrição do Backbone}

O \textbf{EfficientNetB1} é a primeira escala acima do B0 na família EfficientNet,
com resolução nativa de $240\times240$ e maior número de parâmetros. Mantém a mesma
estratégia de \emph{compound scaling}.

\subsection{Hiperparâmetros}

\begin{parametros}
Backbone              & EfficientNetB1 (ImageNet, congelado) \\
Taxa de Aprendizado   & $1 \times 10^{-4}$ \\
Imagem de Entrada     & $240 \times 240 \times 3$ \\
Tamanho do Batch      & 32 \\
Épocas Máximas        & 30 \\
Otimizador            & Adam \\
Pré-processamento     & \texttt{efficientnet.preprocess\_input} \\
\end{parametros}

\subsection{Resultados por Rodada}

\tabelarodadas{Resultados por rodada --- EfficientNetB1}%
              {../results/latex/tabela_rodadas_efficientnetb1.tex}

\subsection{Estatísticas Resumo}

\tabelaresumo{Resumo estatístico --- EfficientNetB1}%
             {../results/latex/tabela_resumo_efficientnetb1.tex}

\subsection{Gráficos}

\begin{figure}[H]
  \centering
  \includegraphics[width=0.85\textwidth]{efficientnetb1_01_acuracia_por_rodada.pdf}
  \caption{Acurácia por rodada --- EfficientNetB1}
\end{figure}

\begin{figure}[H]
  \centering
  \includegraphics[width=0.85\textwidth]{efficientnetb1_02_boxplot_metricas.pdf}
  \caption{Boxplot das métricas --- EfficientNetB1}
\end{figure}

\begin{figure}[H]
  \centering
  \includegraphics[width=0.85\textwidth]{efficientnetb1_03_media_mais_desvio.pdf}
  \caption{Média $\pm$ desvio padrão --- EfficientNetB1}
\end{figure}

\subsection{Análise}

O EfficientNetB1 apresentou desempenho idêntico ao B0: acurácia média $0{,}9983$
e desvio de $0{,}0053$. O custo computacional foi ligeiramente maior ($54{,}3$~s
de treinamento em média \emph{vs.} $47{,}1$~s do B0), sem ganho preditivo. Isso
indica que o aumento de complexidade do B1 não agrega valor para esta tarefa binária.

% ==============================================================
\section{Modelo 6 --- EfficientNetB2}
\label{sec:efficientnetb2}
% ==============================================================

\subsection{Descrição do Backbone}

O \textbf{EfficientNetB2} continua o escalonamento da família EfficientNet,
com resolução nativa de $260\times260$. É o maior dos três variantes avaliados
neste trabalho.

\subsection{Hiperparâmetros}

\begin{parametros}
Backbone              & EfficientNetB2 (ImageNet, congelado) \\
Taxa de Aprendizado   & $1 \times 10^{-4}$ \\
Imagem de Entrada     & $260 \times 260 \times 3$ \\
Tamanho do Batch      & 32 \\
Épocas Máximas        & 30 \\
Otimizador            & Adam \\
Pré-processamento     & \texttt{efficientnet.preprocess\_input} \\
\end{parametros}

\subsection{Resultados por Rodada}

\tabelarodadas{Resultados por rodada --- EfficientNetB2}%
              {../results/latex/tabela_rodadas_efficientnetb2.tex}

\subsection{Estatísticas Resumo}

\tabelaresumo{Resumo estatístico --- EfficientNetB2}%
             {../results/latex/tabela_resumo_efficientnetb2.tex}

\subsection{Gráficos}

\begin{figure}[H]
  \centering
  \includegraphics[width=0.85\textwidth]{efficientnetb2_01_acuracia_por_rodada.pdf}
  \caption{Acurácia por rodada --- EfficientNetB2}
\end{figure}

\begin{figure}[H]
  \centering
  \includegraphics[width=0.85\textwidth]{efficientnetb2_02_boxplot_metricas.pdf}
  \caption{Boxplot das métricas --- EfficientNetB2}
\end{figure}

\begin{figure}[H]
  \centering
  \includegraphics[width=0.85\textwidth]{efficientnetb2_03_media_mais_desvio.pdf}
  \caption{Média $\pm$ desvio padrão --- EfficientNetB2}
\end{figure}

\subsection{Análise}

O EfficientNetB2 é um dos dois modelos com \textbf{desempenho perfeito}: acurácia,
precisão, recall e F1-score médios iguais a $1{,}0000$ com desvio padrão $0{,}0000$
em todas as 10~rodadas. O treinamento demora em média $55{,}7$~s, tornando-o mais
lento que o B0 porém com resultado absolutamente estável.

% ==============================================================
\section{Modelo 7 --- ResNet50}
\label{sec:resnet50}
% ==============================================================

\subsection{Descrição do Backbone}

A \textbf{ResNet50} \cite{he2016deep} é uma rede residual com 50 camadas e conexões
de atalho (\emph{skip connections}) que resolvem o problema do gradiente desvanecente
em redes profundas. Cada bloco residual aprende a diferença (resíduo) entre a entrada
e a saída desejada.

\subsection{Hiperparâmetros}

\begin{parametros}
Backbone              & ResNet50 (ImageNet, congelado) \\
Taxa de Aprendizado   & $1 \times 10^{-4}$ \\
Imagem de Entrada     & $224 \times 224 \times 3$ \\
Tamanho do Batch      & 32 \\
Épocas Máximas        & 30 \\
Otimizador            & Adam \\
Pré-processamento     & \texttt{resnet.preprocess\_input} \\
\end{parametros}

\subsection{Resultados por Rodada}

\tabelarodadas{Resultados por rodada --- ResNet50}%
              {../results/latex/tabela_rodadas_resnet50.tex}

\subsection{Estatísticas Resumo}

\tabelaresumo{Resumo estatístico --- ResNet50}%
             {../results/latex/tabela_resumo_resnet50.tex}

\subsection{Gráficos}

\begin{figure}[H]
  \centering
  \includegraphics[width=0.85\textwidth]{resnet50_01_acuracia_por_rodada.pdf}
  \caption{Acurácia por rodada --- ResNet50}
\end{figure}

\begin{figure}[H]
  \centering
  \includegraphics[width=0.85\textwidth]{resnet50_02_boxplot_metricas.pdf}
  \caption{Boxplot das métricas --- ResNet50}
\end{figure}

\begin{figure}[H]
  \centering
  \includegraphics[width=0.85\textwidth]{resnet50_03_media_mais_desvio.pdf}
  \caption{Média $\pm$ desvio padrão --- ResNet50}
\end{figure}

\subsection{Análise}

A ResNet50 obteve acurácia média de $0{,}9600$ com desvio de $0{,}0387$ (CV = 4,0\%).
Embora seja um dos modelos de TL com menor desempenho médio neste experimento,
permanece bem acima das CNNs do zero. A variabilidade moderada (IC~95\%:
$[0{,}9324;\, 0{,}9876]$) sugere sensibilidade à inicialização da cabeça de
classificação. O tempo de treinamento ($48{,}4$~s) é similar ao EfficientNetB0.

% ==============================================================
\section{Modelo 8 --- DenseNet121}
\label{sec:densenet121}
% ==============================================================

\subsection{Descrição do Backbone}

A \textbf{DenseNet121} \cite{huang2017densely} conecta cada camada a todas as
camadas subsequentes (\emph{dense connections}), promovendo reutilização de
características e fluxo eficiente de gradientes. A DenseNet121 possui
121~camadas, organizadas em 4 blocos densos com camadas de transição entre eles.

\subsection{Hiperparâmetros}

\begin{parametros}
Backbone              & DenseNet121 (ImageNet, congelado) \\
Taxa de Aprendizado   & $1 \times 10^{-4}$ \\
Imagem de Entrada     & $224 \times 224 \times 3$ \\
Tamanho do Batch      & 32 \\
Épocas Máximas        & 30 \\
Otimizador            & Adam \\
Pré-processamento     & \texttt{densenet.preprocess\_input} \\
\end{parametros}

\subsection{Resultados por Rodada}

\tabelarodadas{Resultados por rodada --- DenseNet121}%
              {../results/latex/tabela_rodadas_densenet121.tex}

\subsection{Estatísticas Resumo}

\tabelaresumo{Resumo estatístico --- DenseNet121}%
             {../results/latex/tabela_resumo_densenet121.tex}

\subsection{Gráficos}

\begin{figure}[H]
  \centering
  \includegraphics[width=0.85\textwidth]{densenet121_01_acuracia_por_rodada.pdf}
  \caption{Acurácia por rodada --- DenseNet121}
\end{figure}

\begin{figure}[H]
  \centering
  \includegraphics[width=0.85\textwidth]{densenet121_02_boxplot_metricas.pdf}
  \caption{Boxplot das métricas --- DenseNet121}
\end{figure}

\begin{figure}[H]
  \centering
  \includegraphics[width=0.85\textwidth]{densenet121_03_media_mais_desvio.pdf}
  \caption{Média $\pm$ desvio padrão --- DenseNet121}
\end{figure}

\subsection{Análise}

A DenseNet121 alcançou acurácia média de $0{,}9833$ com desvio de $0{,}0136$
(CV = 1,4\%). O IC~95\% é $[0{,}9736;\, 0{,}9931]$, indicando alta estabilidade.
O tempo médio de treinamento foi de $57{,}2$~s, o que a torna um dos modelos
mais lentos nessa resolução, reflexo do alto número de conexões densas internas.

% ==============================================================
\section{Modelo 9 --- MobileNetV3Small}
\label{sec:mobilenetv3small}
% ==============================================================

\subsection{Descrição do Backbone}

O \textbf{MobileNetV3Small} \cite{howard2019searching} é uma versão compacta da
família MobileNetV3, projetada para dispositivos de recursos muito limitados.
Incorpora blocos \emph{squeeze-and-excitation} e a função de ativação \emph{hard-swish}.
É a variante menor entre as duas versões do V3 avaliadas.

\subsection{Hiperparâmetros}

\begin{parametros}
Backbone              & MobileNetV3Small (ImageNet, congelado) \\
Taxa de Aprendizado   & $1 \times 10^{-4}$ \\
Imagem de Entrada     & $224 \times 224 \times 3$ \\
Tamanho do Batch      & 32 \\
Épocas Máximas        & 30 \\
Otimizador            & Adam \\
Pré-processamento     & \texttt{mobilenet\_v3.preprocess\_input} \\
\end{parametros}

\subsection{Resultados por Rodada}

\tabelarodadas{Resultados por rodada --- MobileNetV3Small}%
              {../results/latex/tabela_rodadas_mobilenetv3small.tex}

\subsection{Estatísticas Resumo}

\tabelaresumo{Resumo estatístico --- MobileNetV3Small}%
             {../results/latex/tabela_resumo_mobilenetv3small.tex}

\subsection{Gráficos}

\begin{figure}[H]
  \centering
  \includegraphics[width=0.85\textwidth]{mobilenetv3small_01_acuracia_por_rodada.pdf}
  \caption{Acurácia por rodada --- MobileNetV3Small}
\end{figure}

\begin{figure}[H]
  \centering
  \includegraphics[width=0.85\textwidth]{mobilenetv3small_02_boxplot_metricas.pdf}
  \caption{Boxplot das métricas --- MobileNetV3Small}
\end{figure}

\begin{figure}[H]
  \centering
  \includegraphics[width=0.85\textwidth]{mobilenetv3small_03_media_mais_desvio.pdf}
  \caption{Média $\pm$ desvio padrão --- MobileNetV3Small}
\end{figure}

\subsection{Análise}

O MobileNetV3Small atingiu acurácia média de $0{,}9650$ com desvio de $0{,}0631$
(CV = 6,5\%), sendo o modelo de TL com maior variabilidade neste experimento.
O IC~95\% de $[0{,}9199;\, 1{,}0101]$ ultrapassa 1,0 no limite superior (artefato
estatístico aceitável). Apesar de ser o modelo mais compacto e rápido do grupo V3,
sua menor capacidade representacional resulta em instabilidade para esta tarefa.

% ==============================================================
\section{Modelo 10 --- MobileNetV3Large}
\label{sec:mobilenetv3large}
% ==============================================================

\subsection{Descrição do Backbone}

O \textbf{MobileNetV3Large} é a versão de maior capacidade da família MobileNetV3,
com mais filtros e estágios que o \emph{Small}. Mantém a eficiência para hardware
limitado, mas com melhor desempenho em tarefas mais complexas.

\subsection{Hiperparâmetros}

\begin{parametros}
Backbone              & MobileNetV3Large (ImageNet, congelado) \\
Taxa de Aprendizado   & $1 \times 10^{-4}$ \\
Imagem de Entrada     & $224 \times 224 \times 3$ \\
Tamanho do Batch      & 32 \\
Épocas Máximas        & 30 \\
Otimizador            & Adam \\
Pré-processamento     & \texttt{mobilenet\_v3.preprocess\_input} \\
\end{parametros}

\subsection{Resultados por Rodada}

\tabelarodadas{Resultados por rodada --- MobileNetV3Large}%
              {../results/latex/tabela_rodadas_mobilenetv3large.tex}

\subsection{Estatísticas Resumo}

\tabelaresumo{Resumo estatístico --- MobileNetV3Large}%
             {../results/latex/tabela_resumo_mobilenetv3large.tex}

\subsection{Gráficos}

\begin{figure}[H]
  \centering
  \includegraphics[width=0.85\textwidth]{mobilenetv3large_01_acuracia_por_rodada.pdf}
  \caption{Acurácia por rodada --- MobileNetV3Large}
\end{figure}

\begin{figure}[H]
  \centering
  \includegraphics[width=0.85\textwidth]{mobilenetv3large_02_boxplot_metricas.pdf}
  \caption{Boxplot das métricas --- MobileNetV3Large}
\end{figure}

\begin{figure}[H]
  \centering
  \includegraphics[width=0.85\textwidth]{mobilenetv3large_03_media_mais_desvio.pdf}
  \caption{Média $\pm$ desvio padrão --- MobileNetV3Large}
\end{figure}

\subsection{Análise}

O MobileNetV3Large melhorou em relação ao Small: acurácia média de $0{,}9700$
com desvio de $0{,}0312$ (CV = 3,2\%). O IC~95\% é $[0{,}9477;\, 0{,}9923]$.
O tempo de treinamento ($42{,}7$~s) é similar ao MobileNetV2. A versão Large
do V3 mostra melhor equilíbrio entre estabilidade e eficiência comparado à versão Small.

% ==============================================================
\section{Modelo 11 --- NASNetMobile}
\label{sec:nasnetmobile}
% ==============================================================

\subsection{Descrição do Backbone}

O \textbf{NASNetMobile} \cite{zoph2018learning} é uma arquitetura obtida por
\emph{Neural Architecture Search} (NAS), em que o design da rede é automaticamente
otimizado por um algoritmo de busca. A versão \emph{Mobile} é compacta, balanceando
desempenho e eficiência.

\subsection{Hiperparâmetros}

\begin{parametros}
Backbone              & NASNetMobile (ImageNet, congelado) \\
Taxa de Aprendizado   & $1 \times 10^{-4}$ \\
Imagem de Entrada     & $224 \times 224 \times 3$ \\
Tamanho do Batch      & 32 \\
Épocas Máximas        & 30 \\
Otimizador            & Adam \\
Pré-processamento     & \texttt{nasnet.preprocess\_input} \\
\end{parametros}

\subsection{Resultados por Rodada}

\tabelarodadas{Resultados por rodada --- NASNetMobile}%
              {../results/latex/tabela_rodadas_nasnetmobile.tex}

\subsection{Estatísticas Resumo}

\tabelaresumo{Resumo estatístico --- NASNetMobile}%
             {../results/latex/tabela_resumo_nasnetmobile.tex}

\subsection{Gráficos}

\begin{figure}[H]
  \centering
  \includegraphics[width=0.85\textwidth]{nasnetmobile_01_acuracia_por_rodada.pdf}
  \caption{Acurácia por rodada --- NASNetMobile}
\end{figure}

\begin{figure}[H]
  \centering
  \includegraphics[width=0.85\textwidth]{nasnetmobile_02_boxplot_metricas.pdf}
  \caption{Boxplot das métricas --- NASNetMobile}
\end{figure}

\begin{figure}[H]
  \centering
  \includegraphics[width=0.85\textwidth]{nasnetmobile_03_media_mais_desvio.pdf}
  \caption{Média $\pm$ desvio padrão --- NASNetMobile}
\end{figure}

\subsection{Análise}

O NASNetMobile obteve acurácia média de $0{,}9683$ com desvio de $0{,}0254$
(CV = 2,6\%), com IC~95\% em $[0{,}9502;\, 0{,}9865]$. Destaca-se negativamente
pelo maior tempo de predição ($15{,}3$~s), consequência da topologia complexa
gerada pelo NAS, que aumenta a latência de inferência. Seu tempo total médio
($82{,}4$~s) é um dos mais elevados entre os modelos leves.

% ==============================================================
\section{Modelo 12 --- InceptionV3}
\label{sec:inceptionv3}
% ==============================================================

\subsection{Descrição do Backbone}

A \textbf{InceptionV3} \cite{szegedy2016rethinking} utiliza módulos Inception
que paralelizam convoluções de diferentes tamanhos ($1\times1$, $3\times3$,
$5\times5$) e as concatenam. O V3 introduz fatorização de convoluções e normalização
de batch, exigindo entrada mínima de $299\times299$.

\subsection{Hiperparâmetros}

\begin{parametros}
Backbone              & InceptionV3 (ImageNet, congelado) \\
Taxa de Aprendizado   & $1 \times 10^{-4}$ \\
Imagem de Entrada     & $299 \times 299 \times 3$ \\
Tamanho do Batch      & 16 \\
Épocas Máximas        & 30 \\
Otimizador            & Adam \\
Pré-processamento     & \texttt{inception\_v3.preprocess\_input} \\
\end{parametros}

\subsection{Resultados por Rodada}

\tabelarodadas{Resultados por rodada --- InceptionV3}%
              {../results/latex/tabela_rodadas_inceptionv3.tex}

\subsection{Estatísticas Resumo}

\tabelaresumo{Resumo estatístico --- InceptionV3}%
             {../results/latex/tabela_resumo_inceptionv3.tex}

\subsection{Gráficos}

\begin{figure}[H]
  \centering
  \includegraphics[width=0.85\textwidth]{inceptionv3_01_acuracia_por_rodada.pdf}
  \caption{Acurácia por rodada --- InceptionV3}
\end{figure}

\begin{figure}[H]
  \centering
  \includegraphics[width=0.85\textwidth]{inceptionv3_02_boxplot_metricas.pdf}
  \caption{Boxplot das métricas --- InceptionV3}
\end{figure}

\begin{figure}[H]
  \centering
  \includegraphics[width=0.85\textwidth]{inceptionv3_03_media_mais_desvio.pdf}
  \caption{Média $\pm$ desvio padrão --- InceptionV3}
\end{figure}

\subsection{Análise}

A InceptionV3 obteve acurácia média de $0{,}9833$ com desvio de $0{,}0111$
(CV = 1,1\%), com IC~95\% em $[0{,}9754;\, 0{,}9913]$. Apresenta alto desempenho
preditivo, porém com o maior tempo de treinamento entre os modelos de resolução
$299\times299$, chegando a $99{,}2$~s em média. O batch de 16 (necessário pela
resolução maior) contribui para esse custo.

% ==============================================================
\section{Modelo 13 --- Xception}
\label{sec:xception}
% ==============================================================

\subsection{Descrição do Backbone}

A \textbf{Xception} \cite{chollet2017xception} é uma extensão da Inception que
substitui os módulos Inception por \emph{depthwise separable convolutions} em
todas as camadas. O nome vem de \emph{Extreme Inception}. Assim como a InceptionV3,
requer entrada de $299\times299$.

\subsection{Hiperparâmetros}

\begin{parametros}
Backbone              & Xception (ImageNet, congelado) \\
Taxa de Aprendizado   & $1 \times 10^{-4}$ \\
Imagem de Entrada     & $299 \times 299 \times 3$ \\
Tamanho do Batch      & 16 \\
Épocas Máximas        & 30 \\
Otimizador            & Adam \\
Pré-processamento     & \texttt{xception.preprocess\_input} \\
\end{parametros}

\subsection{Resultados por Rodada}

\tabelarodadas{Resultados por rodada --- Xception}%
              {../results/latex/tabela_rodadas_xception.tex}

\subsection{Estatísticas Resumo}

\tabelaresumo{Resumo estatístico --- Xception}%
             {../results/latex/tabela_resumo_xception.tex}

\subsection{Gráficos}

\begin{figure}[H]
  \centering
  \includegraphics[width=0.85\textwidth]{xception_01_acuracia_por_rodada.pdf}
  \caption{Acurácia por rodada --- Xception}
\end{figure}

\begin{figure}[H]
  \centering
  \includegraphics[width=0.85\textwidth]{xception_02_boxplot_metricas.pdf}
  \caption{Boxplot das métricas --- Xception}
\end{figure}

\begin{figure}[H]
  \centering
  \includegraphics[width=0.85\textwidth]{xception_03_media_mais_desvio.pdf}
  \caption{Média $\pm$ desvio padrão --- Xception}
\end{figure}

\subsection{Análise}

A Xception é o segundo modelo com \textbf{desempenho perfeito} em todas as rodadas:
acurácia, precisão, recall e F1-score iguais a $1{,}0000$ com desvio $0{,}0000$.
O custo é o maior de todos os modelos: $110{,}7$~s de treinamento e $116{,}4$~s
no total. Para esta tarefa binária, a Xception alcança o limite superior do desempenho,
mas com o maior tempo de processamento.

% ==============================================================
\section{Análise Comparativa dos 13 Modelos}
% ==============================================================

\subsection{Tabela Comparativa de Médias --- Métricas de Classificação}

\begin{table}[H]
\centering
\caption{Média das métricas de classificação por modelo (10 rodadas)}
\begin{adjustbox}{max width=\textwidth, center}
\begin{tabular}{lrrrr}
\toprule
Metrica & Acuracia & F1_score & Precision & Recall \\
Modelo &  &  &  &  \\
\midrule
cnn_from_scratch & 0.7400 & 0.7267 & 0.7882 & 0.7400 \\
densenet121 & 0.9833 & 0.9833 & 0.9842 & 0.9833 \\
efficientnetb0 & 0.9983 & 0.9983 & 0.9984 & 0.9983 \\
efficientnetb1 & 0.9983 & 0.9983 & 0.9984 & 0.9983 \\
efficientnetb2 & 1.0000 & 1.0000 & 1.0000 & 1.0000 \\
inceptionv3 & 0.9833 & 0.9833 & 0.9841 & 0.9833 \\
mobilenetv2 & 0.9700 & 0.9699 & 0.9723 & 0.9700 \\
mobilenetv3large & 0.9700 & 0.9699 & 0.9717 & 0.9700 \\
mobilenetv3small & 0.9650 & 0.9641 & 0.9721 & 0.9650 \\
nasnetmobile & 0.9683 & 0.9683 & 0.9700 & 0.9683 \\
resnet50 & 0.9600 & 0.9597 & 0.9650 & 0.9600 \\
xception & 1.0000 & 1.0000 & 1.0000 & 1.0000 \\
\bottomrule
\end{tabular}

\end{adjustbox}
\end{table}

\subsection{Tabela Comparativa de Médias --- Tempos de Execução}

\begin{table}[H]
\centering
\caption{Média dos tempos de execução por modelo, em segundos (10 rodadas)}
\begin{adjustbox}{max width=\textwidth, center}
\begin{tabular}{lrrr}
\toprule
Metrica & Tempo_Pred_s & Tempo_Total_s & Tempo_Treino_s \\
Modelo &  &  &  \\
\midrule
cnn_from_scratch & 6.5507 & 40.9079 & 34.3572 \\
densenet121 & 11.3072 & 68.5330 & 57.2258 \\
efficientnetb0 & 7.8116 & 54.9136 & 47.1019 \\
efficientnetb1 & 9.7962 & 64.0731 & 54.2768 \\
efficientnetb2 & 9.9932 & 65.7119 & 55.7187 \\
inceptionv3 & 8.0358 & 107.2421 & 99.2064 \\
mobilenetv2 & 5.8863 & 46.1598 & 40.2735 \\
mobilenetv3large & 6.5665 & 49.2361 & 42.6697 \\
mobilenetv3small & 5.8935 & 46.0466 & 40.1531 \\
nasnetmobile & 15.2543 & 82.3962 & 67.1419 \\
resnet50 & 7.9943 & 56.3756 & 48.3813 \\
xception & 5.6374 & 116.3571 & 110.7196 \\
\bottomrule
\end{tabular}

\end{adjustbox}
\end{table}

\subsection{Tabela Resumo Consolidada dos 13 Modelos}

A tabela a seguir apresenta todos os resumos estatísticos de todos os 13~modelos
concatenados, incluindo média, desvio padrão, coeficiente de variação e
intervalo de confiança de 95\% para cada métrica.

\begin{table}[H]
\centering
\caption{Resumo estatístico consolidado dos 13 modelos}
\begin{adjustbox}{max width=\textwidth, center}
\begin{tabular}{lllllll}
\toprule
Metrica & Media & Desvio_Padrao & Coef_Variacao & IC95_Lower & IC95_Upper & Modelo \\
\midrule
Acuracia & 0.7400 & 0.0903 & 0.1221 & 0.6754 & 0.8046 & cnn_from_scratch \\
Precision & 0.7882 & 0.0846 & 0.1074 & 0.7277 & 0.8488 & cnn_from_scratch \\
Recall & 0.7400 & 0.0903 & 0.1221 & 0.6754 & 0.8046 & cnn_from_scratch \\
F1_score & 0.7267 & 0.0986 & 0.1356 & 0.6562 & 0.7972 & cnn_from_scratch \\
Tempo_Treino_s & 34.3572 & 15.0289 & 0.4374 & 23.6069 & 45.1075 & cnn_from_scratch \\
Tempo_Pred_s & 6.5507 & 9.3739 & 1.4310 & -0.1545 & 13.2559 & cnn_from_scratch \\
Tempo_Total_s & 40.9079 & 23.6063 & 0.5771 & 24.0222 & 57.7937 & cnn_from_scratch \\
Acuracia & 0.9700 & 0.0205 & 0.0211 & 0.9553 & 0.9847 & mobilenetv2 \\
Precision & 0.9723 & 0.0180 & 0.0185 & 0.9594 & 0.9852 & mobilenetv2 \\
Recall & 0.9700 & 0.0205 & 0.0211 & 0.9553 & 0.9847 & mobilenetv2 \\
F1_score & 0.9699 & 0.0206 & 0.0212 & 0.9552 & 0.9847 & mobilenetv2 \\
Tempo_Treino_s & 40.2735 & 3.0055 & 0.0746 & 38.1236 & 42.4234 & mobilenetv2 \\
Tempo_Pred_s & 5.8863 & 0.4626 & 0.0786 & 5.5554 & 6.2172 & mobilenetv2 \\
Tempo_Total_s & 46.1598 & 2.9721 & 0.0644 & 44.0338 & 48.2857 & mobilenetv2 \\
Acuracia & 0.9983 & 0.0053 & 0.0053 & 0.9946 & 1.0021 & efficientnetb0 \\
Precision & 0.9984 & 0.0051 & 0.0051 & 0.9947 & 1.0020 & efficientnetb0 \\
Recall & 0.9983 & 0.0053 & 0.0053 & 0.9946 & 1.0021 & efficientnetb0 \\
F1_score & 0.9983 & 0.0053 & 0.0053 & 0.9946 & 1.0021 & efficientnetb0 \\
Tempo_Treino_s & 47.1019 & 1.0835 & 0.0230 & 46.3269 & 47.8770 & efficientnetb0 \\
Tempo_Pred_s & 7.8116 & 0.4903 & 0.0628 & 7.4609 & 8.1623 & efficientnetb0 \\
Tempo_Total_s & 54.9136 & 0.9837 & 0.0179 & 54.2099 & 55.6172 & efficientnetb0 \\
Acuracia & 0.9983 & 0.0053 & 0.0053 & 0.9946 & 1.0021 & efficientnetb1 \\
Precision & 0.9984 & 0.0051 & 0.0051 & 0.9947 & 1.0020 & efficientnetb1 \\
Recall & 0.9983 & 0.0053 & 0.0053 & 0.9946 & 1.0021 & efficientnetb1 \\
F1_score & 0.9983 & 0.0053 & 0.0053 & 0.9946 & 1.0021 & efficientnetb1 \\
Tempo_Treino_s & 54.2768 & 0.7713 & 0.0142 & 53.7251 & 54.8285 & efficientnetb1 \\
Tempo_Pred_s & 9.7962 & 0.3557 & 0.0363 & 9.5418 & 10.0506 & efficientnetb1 \\
Tempo_Total_s & 64.0731 & 0.6945 & 0.0108 & 63.5763 & 64.5698 & efficientnetb1 \\
Acuracia & 1.0000 & 0.0000 & 0.0000 & 1.0000 & 1.0000 & efficientnetb2 \\
Precision & 1.0000 & 0.0000 & 0.0000 & 1.0000 & 1.0000 & efficientnetb2 \\
Recall & 1.0000 & 0.0000 & 0.0000 & 1.0000 & 1.0000 & efficientnetb2 \\
F1_score & 1.0000 & 0.0000 & 0.0000 & 1.0000 & 1.0000 & efficientnetb2 \\
Tempo_Treino_s & 55.7187 & 0.9358 & 0.0168 & 55.0493 & 56.3881 & efficientnetb2 \\
Tempo_Pred_s & 9.9932 & 0.3444 & 0.0345 & 9.7469 & 10.2395 & efficientnetb2 \\
Tempo_Total_s & 65.7119 & 0.8992 & 0.0137 & 65.0687 & 66.3551 & efficientnetb2 \\
Acuracia & 0.9600 & 0.0387 & 0.0403 & 0.9324 & 0.9876 & resnet50 \\
Precision & 0.9650 & 0.0323 & 0.0335 & 0.9419 & 0.9882 & resnet50 \\
Recall & 0.9600 & 0.0387 & 0.0403 & 0.9324 & 0.9876 & resnet50 \\
F1_score & 0.9597 & 0.0390 & 0.0407 & 0.9318 & 0.9877 & resnet50 \\
Tempo_Treino_s & 48.3813 & 2.6453 & 0.0547 & 46.4891 & 50.2736 & resnet50 \\
Tempo_Pred_s & 7.9943 & 1.4309 & 0.1790 & 6.9707 & 9.0178 & resnet50 \\
Tempo_Total_s & 56.3756 & 2.8893 & 0.0513 & 54.3089 & 58.4423 & resnet50 \\
Acuracia & 0.9833 & 0.0136 & 0.0138 & 0.9736 & 0.9931 & densenet121 \\
Precision & 0.9842 & 0.0128 & 0.0130 & 0.9750 & 0.9933 & densenet121 \\
Recall & 0.9833 & 0.0136 & 0.0138 & 0.9736 & 0.9931 & densenet121 \\
F1_score & 0.9833 & 0.0136 & 0.0139 & 0.9736 & 0.9931 & densenet121 \\
Tempo_Treino_s & 57.2258 & 1.7831 & 0.0312 & 55.9503 & 58.5012 & densenet121 \\
Tempo_Pred_s & 11.3072 & 0.2795 & 0.0247 & 11.1073 & 11.5072 & densenet121 \\
Tempo_Total_s & 68.5330 & 1.8383 & 0.0268 & 67.2180 & 69.8480 & densenet121 \\
Acuracia & 0.9650 & 0.0631 & 0.0654 & 0.9199 & 1.0101 & mobilenetv3small \\
Precision & 0.9721 & 0.0457 & 0.0470 & 0.9394 & 1.0048 & mobilenetv3small \\
Recall & 0.9650 & 0.0631 & 0.0654 & 0.9199 & 1.0101 & mobilenetv3small \\
F1_score & 0.9641 & 0.0656 & 0.0680 & 0.9172 & 1.0110 & mobilenetv3small \\
Tempo_Treino_s & 40.1531 & 0.6306 & 0.0157 & 39.7021 & 40.6042 & mobilenetv3small \\
Tempo_Pred_s & 5.8935 & 0.4665 & 0.0792 & 5.5598 & 6.2271 & mobilenetv3small \\
Tempo_Total_s & 46.0466 & 0.8977 & 0.0195 & 45.4045 & 46.6887 & mobilenetv3small \\
Acuracia & 0.9700 & 0.0312 & 0.0322 & 0.9477 & 0.9923 & mobilenetv3large \\
Precision & 0.9717 & 0.0285 & 0.0294 & 0.9513 & 0.9921 & mobilenetv3large \\
Recall & 0.9700 & 0.0312 & 0.0322 & 0.9477 & 0.9923 & mobilenetv3large \\
F1_score & 0.9699 & 0.0314 & 0.0324 & 0.9475 & 0.9924 & mobilenetv3large \\
Tempo_Treino_s & 42.6697 & 0.7924 & 0.0186 & 42.1029 & 43.2365 & mobilenetv3large \\
Tempo_Pred_s & 6.5665 & 0.4238 & 0.0645 & 6.2633 & 6.8696 & mobilenetv3large \\
Tempo_Total_s & 49.2361 & 0.7459 & 0.0151 & 48.7026 & 49.7697 & mobilenetv3large \\
Acuracia & 0.9683 & 0.0254 & 0.0262 & 0.9502 & 0.9865 & nasnetmobile \\
Precision & 0.9700 & 0.0237 & 0.0244 & 0.9530 & 0.9869 & nasnetmobile \\
Recall & 0.9683 & 0.0254 & 0.0262 & 0.9502 & 0.9865 & nasnetmobile \\
F1_score & 0.9683 & 0.0255 & 0.0263 & 0.9501 & 0.9865 & nasnetmobile \\
Tempo_Treino_s & 67.1419 & 2.6777 & 0.0399 & 65.2265 & 69.0573 & nasnetmobile \\
Tempo_Pred_s & 15.2543 & 0.2042 & 0.0134 & 15.1082 & 15.4004 & nasnetmobile \\
Tempo_Total_s & 82.3962 & 2.6775 & 0.0325 & 80.4810 & 84.3114 & nasnetmobile \\
Acuracia & 0.9833 & 0.0111 & 0.0113 & 0.9754 & 0.9913 & inceptionv3 \\
Precision & 0.9841 & 0.0104 & 0.0106 & 0.9766 & 0.9915 & inceptionv3 \\
Recall & 0.9833 & 0.0111 & 0.0113 & 0.9754 & 0.9913 & inceptionv3 \\
F1_score & 0.9833 & 0.0111 & 0.0113 & 0.9754 & 0.9913 & inceptionv3 \\
Tempo_Treino_s & 99.2064 & 9.4282 & 0.0950 & 92.4623 & 105.9504 & inceptionv3 \\
Tempo_Pred_s & 8.0358 & 0.4683 & 0.0583 & 7.7008 & 8.3707 & inceptionv3 \\
Tempo_Total_s & 107.2421 & 9.2435 & 0.0862 & 100.6302 & 113.8541 & inceptionv3 \\
Acuracia & 1.0000 & 0.0000 & 0.0000 & 1.0000 & 1.0000 & xception \\
Precision & 1.0000 & 0.0000 & 0.0000 & 1.0000 & 1.0000 & xception \\
Recall & 1.0000 & 0.0000 & 0.0000 & 1.0000 & 1.0000 & xception \\
F1_score & 1.0000 & 0.0000 & 0.0000 & 1.0000 & 1.0000 & xception \\
Tempo_Treino_s & 110.7196 & 1.8956 & 0.0171 & 109.3637 & 112.0755 & xception \\
Tempo_Pred_s & 5.6374 & 0.4644 & 0.0824 & 5.3052 & 5.9697 & xception \\
Tempo_Total_s & 116.3571 & 2.0847 & 0.0179 & 114.8658 & 117.8483 & xception \\
\bottomrule
\end{tabular}

\end{adjustbox}
\end{table}

\subsection{Discussão Geral}

\subsubsection{Modelos de Melhor Desempenho}

Os modelos \textbf{EfficientNetB2} e \textbf{Xception} atingiram acurácia perfeita
($1{,}0000$) em todas as 10~rodadas, com desvio padrão zero. Eles representam
a fronteira superior de desempenho para esta tarefa binária.

\subsubsection{Família EfficientNet}

A família EfficientNet demonstrou excelente desempenho em todos os três variantes:
\begin{itemize}
  \item \textbf{B0}: acurácia $0{,}9983$ --- melhor custo-benefício da família
  \item \textbf{B1}: acurácia $0{,}9983$ --- sem ganho em relação ao B0, maior custo
  \item \textbf{B2}: acurácia $1{,}0000$ --- melhor desempenho, custo moderado
\end{itemize}

\subsubsection{Modelos Leves (\textit{Mobile})}

MobileNetV2, MobileNetV3Large e NASNetMobile apresentam bom desempenho
(acurácia $\geq 0{,}97$) com tempos de treinamento reduzidos ($\approx 40$--50~s),
sendo adequados para cenários com restrições computacionais.

\subsubsection{Impacto da Arquitetura CNN do Zero}

A comparação entre as CNNs treinadas do zero e os modelos de TL ilustra
claramente o \textbf{poder do Transfer Learning}: todos os 11~backbones
pré-treinados superaram as CNNs do zero, que atingiram no máximo $0{,}74$ de acurácia.
Isso confirma que os pesos aprendidos no ImageNet transferem representações visuais
úteis mesmo para domínios específicos como moda.

\subsubsection{Relação Desempenho $\times$ Tempo}

\begin{table}[H]
\centering
\caption{Classificação dos modelos por desempenho e custo computacional}
\small
\begin{tabular}{lccc}
\toprule
\textbf{Modelo} & \textbf{Acurácia Média} & \textbf{Tempo Treino (s)} & \textbf{Classificação} \\
\midrule
EfficientNetB2  & 1.0000 & 55.7 & Melhor desempenho \\
Xception        & 1.0000 & 110.7 & Melhor desempenho, mais lento \\
EfficientNetB0  & 0.9983 & 47.1 & Ótimo custo-benefício \\
EfficientNetB1  & 0.9983 & 54.3 & Similar ao B0, mais lento \\
DenseNet121     & 0.9833 & 57.2 & Alto desempenho \\
InceptionV3     & 0.9833 & 99.2 & Alto desempenho, lento \\
MobileNetV2     & 0.9700 & 40.3 & Bom, muito rápido \\
MobileNetV3L    & 0.9700 & 42.7 & Bom, rápido \\
NASNetMobile    & 0.9683 & 67.1 & Bom, lento na inferência \\
MobileNetV3S    & 0.9650 & 40.2 & Aceitável, mais variável \\
ResNet50        & 0.9600 & 48.4 & Bom, CV moderado \\
CNN (lr=5e-4)   & 0.7400 & 34.4 & Fraco, muito variável \\
CNN (lr=1e-4)   & 0.6567 & 36.9 & Fraco, muito variável \\
\bottomrule
\end{tabular}
\end{table}

% ==============================================================
\section{Conclusões}
% ==============================================================

Este trabalho avaliou sistematicamente 13 modelos de CNNs para a classificação binária
de imagens de moda (bolsas \emph{vs.} calças), utilizando o protocolo Repeated Holdout
com 10 sementes, garantindo a replicabilidade e robustez estatística dos resultados.

Os principais achados são:

\begin{enumerate}
  \item \textbf{Transfer Learning supera CNNs do zero}: todos os 11 backbones
        pré-treinados no ImageNet superaram significativamente as CNNs treinadas
        do zero, confirmando a eficácia do TL mesmo para domínios específicos.

  \item \textbf{EfficientNetB2 e Xception}: atingiram desempenho perfeito em todas
        as rodadas. O B2 é preferível em custo-benefício pois é mais rápido.

  \item \textbf{EfficientNetB0}: melhor relação desempenho/custo entre os modelos
        avaliados, com acurácia de $0{,}9983$ e treinamento em $\approx 47$~s.

  \item \textbf{MobileNetV2}: ideal para aplicações com restrições de tempo,
        atingindo $0{,}9700$ de acurácia em apenas $\approx 40$~s.

  \item \textbf{Taxa de aprendizado nas CNNs do zero}: a taxa $5\times10^{-4}$
        superou $1\times10^{-4}$, mas ambas apresentaram alta variabilidade
        e desempenho muito inferior ao TL.
\end{enumerate}

\textbf{Recomendação}: Para implantação em produção, o \textbf{EfficientNetB2}
é o modelo recomendado, com acurácia perfeita e tempo de treinamento razoável.
Para cenários com restrições de hardware, o \textbf{MobileNetV2} ou o
\textbf{EfficientNetB0} são alternativas eficientes.

% ==============================================================
\section*{Referências}
% ==============================================================
\addcontentsline{toc}{section}{Referências}

\begin{thebibliography}{99}

\bibitem{kingma2014adam}
D. P. Kingma, J. Ba,
``Adam: A Method for Stochastic Optimization,''
\emph{ICLR}, 2015.

\bibitem{sandler2018mobilenetv2}
M. Sandler, A. Howard, M. Zhu, A. Zhmoginov, L.-C. Chen,
``MobileNetV2: Inverted Residuals and Linear Bottlenecks,''
\emph{CVPR}, 2018.

\bibitem{tan2019efficientnet}
M. Tan, Q. V. Le,
``EfficientNet: Rethinking Model Scaling for Convolutional Neural Networks,''
\emph{ICML}, 2019.

\bibitem{he2016deep}
K. He, X. Zhang, S. Ren, J. Sun,
``Deep Residual Learning for Image Recognition,''
\emph{CVPR}, 2016.

\bibitem{huang2017densely}
G. Huang, Z. Liu, L. van der Maaten, K. Q. Weinberger,
``Densely Connected Convolutional Networks,''
\emph{CVPR}, 2017.

\bibitem{howard2019searching}
A. Howard, R. Pang, H. Adam, Q. V. Le, M. Sandler, B. Chen,
W. Wang, L.-C. Chen, M. Tan, G. Chu, V. Vasudevan, Y. Zhu,
``Searching for MobileNetV3,''
\emph{ICCV}, 2019.

\bibitem{zoph2018learning}
B. Zoph, V. Vasudevan, J. Shlens, Q. V. Le,
``Learning Transferable Architectures for Scalable Image Recognition,''
\emph{CVPR}, 2018.

\bibitem{szegedy2016rethinking}
C. Szegedy, V. Vanhoucke, S. Ioffe, J. Shlens, Z. Wojna,
``Rethinking the Inception Architecture for Computer Vision,''
\emph{CVPR}, 2016.

\bibitem{chollet2017xception}
F. Chollet,
``Xception: Deep Learning with Depthwise Separable Convolutions,''
\emph{CVPR}, 2017.

\end{thebibliography}

\end{document}
